\documentclass[12pt]{article}

\usepackage[utf8]{inputenc}
\usepackage[T1]{fontenc}
\usepackage{lmodern}
\usepackage[spanish]{babel}
\usepackage{booktabs}
\usepackage{amsmath}
\usepackage{forest}
\usepackage{float}
\usepackage{listings}
\usepackage{xcolor}
\usepackage{tikz}

\definecolor{codegreen}{rgb}{0,0.6,0}
\definecolor{codegray}{rgb}{0.5,0.5,0.5}
\definecolor{codepurple}{rgb}{0.58,0,0.82}
\definecolor{backcolour}{rgb}{0.95,0.95,0.92}

\lstdefinestyle{mystyle}{
    backgroundcolor=\color{backcolour},   
    commentstyle=\color{codegreen},
    keywordstyle=\color{magenta},
    numberstyle=\tiny\color{codegray},
    stringstyle=\color{codepurple},
    basicstyle=\ttfamily\footnotesize,
    breakatwhitespace=false,         
    breaklines=true,                 
    captionpos=b,                    
    keepspaces=true,                 
    numbers=left,                    
    numbersep=5pt,                  
    showspaces=false,                
    showstringspaces=false,
    showtabs=false,                  
    tabsize=2
}

\lstset{style=mystyle}

\sloppy
\setlength{\parindent}{0pt}

\begin{document}

% Título y materia
\begin{center}
  {\LARGE \textbf{Guía de Ejercicios - Naive Bayes}}\\[0.5em]
  {Analítica de Datos, Universidad de San Andrés}
\end{center}

Si encuentran algún error en el documento o hay alguna duda, mandenmé un mail a rodriguezf@udesa.edu.ar y lo revisamos.

\section{Ejercicios Conceptuales}

\subsection{Verdadero o Falso}

\begin{enumerate}
    \item Es necesario escalar los datos antes de entrenar un modelo de Naive Bayes.
    \item Naive Bayes asume independencia condicional entre las características.
    \item Multinomial Naive Bayes es apropiado para clasificación de texto.
    \item Gaussian Naive Bayes puede manejar características continuas.
    \item Laplace smoothing se usa para evitar probabilidades cero.
    \item Naive Bayes siempre tiene mejor rendimiento que otros algoritmos de clasificación.
    \item La suposición de independencia nunca se cumple en datos reales.
    \item Bernoulli Naive Bayes es útil para características binarias.
    \item Naive Bayes es un algoritmo de aprendizaje discriminativo.
    \item Naive Bayes puede usarse tanto para clasificación como para regresión.
\end{enumerate}

\subsection{Opción Múltiple}

\begin{enumerate}
    \item ¿Por qué se llama ``Naive'' (ingenuo) al algoritmo Naive Bayes?
    \begin{enumerate}
        \item Porque es un algoritmo simple
        \item Porque asume independencia entre características
        \item Porque solo funciona con datos pequeños
        \item Porque no es muy preciso
    \end{enumerate}
    
    \item ¿Cuál es la ventaja principal de Naive Bayes sobre otros algoritmos?
    \begin{enumerate}
        \item Mayor precisión
        \item Manejo automático de missing values
        \item Velocidad de entrenamiento y predicción
        \item No requiere normalización
    \end{enumerate}
    
    \item ¿Qué tipo de Naive Bayes usarías para clasificación de emails spam?
    \begin{enumerate}
        \item Gaussian
        \item Multinomial
        \item Bernoulli
        \item Cualquiera de los anteriores
    \end{enumerate}
    
    \item ¿Cuándo aplicarías Laplace smoothing?
    \begin{enumerate}
        \item Siempre
        \item Solo con datos grandes
        \item Solo con Gaussian Naive Bayes
        \item Cuando hay probabilidades cero
    \end{enumerate}
    
    \item ¿Qué sucede si la suposición de independencia se viola severamente?
    \begin{enumerate}
        \item Puede reducir el rendimiento
        \item El algoritmo no funciona
        \item No hay impacto
        \item Mejora el rendimiento
    \end{enumerate}
\end{enumerate}

\section{Ejercicios Prácticos}

\subsection{Ejercicio 1 - Cálculo Manual de Probabilidades}

El siguiente conjunto de datos brinda información de las condiciones climáticas de 10 días diferentes junto a una etiqueta que nos indica si Juan sale a correr ese día o no. Va a ser nuestro set de entrenamiento para un modelo de clasificación.

\begin{table}[H]
    \centering
    \begin{tabular}{|c|c|c|c|c|}
        \hline
        \textbf{Obs} & \textbf{Lluvia} & \textbf{Temperatura} & \textbf{Viento} & \textbf{¿Entrena?} \\
        \hline
        1 & SI & Frío & Fuerte & NO \\
        \hline
        2 & NO & Templado & Débil & SI \\
        \hline
        3 & SI & Cálido & Fuerte & SI \\
        \hline
        4 & SI & Templado & Fuerte & NO \\
        \hline
        5 & NO & Templado & Débil & SI \\
        \hline
        6 & NO & Cálido & Fuerte & NO \\
        \hline
        7 & NO & Frío & Fuerte & SI \\
        \hline
        8 & SI & Frío & Débil & NO \\
        \hline
        9 & NO & Cálido & Débil & NO \\
        \hline
        10 & SI & Templado & Fuerte & SI \\
        \hline
    \end{tabular}
\end{table}

\begin{enumerate}
    \item[a.] Entrenar un modelo de Naive Bayes
    \item[b.] Considerando la siguiente muestra de test, calcular las probabilidades de entrenar y no entrenar para cada observación, elegir una métrica de bondad de ajuste del modelo y reportar su valor:
\end{enumerate}

\begin{table}[H]
    \centering
    \begin{tabular}{|c|c|c|c|c|}
        \hline
        \textbf{Obs} & \textbf{Lluvia} & \textbf{Temperatura} & \textbf{Viento} & \textbf{¿Entrena?} \\
        \hline
        11 & NO & Templado & Débil & SI \\
        \hline
        12 & SI & Frío & Fuerte & NO \\
        \hline
        13 & SI & Cálido & Débil & SI \\
        \hline
    \end{tabular}
\end{table}

\subsection{Ejercicio 2 - Clasificación de Spam}

Considerar el siguiente dataset para clasificar emails como ``Spam'' o ``No Spam'' basándose en la presencia de ciertas palabras:

\begin{table}[H]
    \centering
    \begin{tabular}{|c|c|c|c|c|}
        \hline
        \textbf{Email} & \textbf{Gratis} & \textbf{Dinero} & \textbf{Oferta} & \textbf{Clase} \\
        \hline
        1 & 1 & 0 & 1 & Spam \\
        \hline
        2 & 0 & 1 & 0 & Spam \\
        \hline
        3 & 1 & 1 & 1 & Spam \\
        \hline
        4 & 0 & 0 & 0 & No Spam \\
        \hline
        5 & 0 & 0 & 1 & No Spam \\
        \hline
        6 & 1 & 0 & 0 & Spam \\
        \hline
    \end{tabular}
\end{table}

\begin{enumerate}
    \item[a.] Calcular las probabilidades a priori de cada clase
    \item[b.] Calcular las probabilidades condicionales para cada característica
    \item[c.] Clasificar un nuevo email con características: Gratis=1, Dinero=0, Oferta=0
    \item[d.] ¿Qué sucede si aplicamos Laplace smoothing con $\alpha=1$? Recalcular la clasificación del inciso c.
\end{enumerate}

\subsection{Ejercicio 3 - Bag of Words y NLP}

Cuando usamos la técnica Bag of Words para preprocesar los datos previo al entrenamiento de un modelo de NLP (Natural Language Processing o Procesamiento de Lenguaje Natural), ¿Qué características tiene el dataset luego de ese preprocesamiento? Supongamos que queremos clasificar artículos periodísticos en cuatro categorías diferentes (Economía, Deportes, Cultura y Ciencia), ¿Cómo funciona un modelo de Naive Bayes en este caso?

\subsection{Ejercicio 4 - Implementación en Python}

Implementar en Python un modelo de Naive Bayes para el dataset del Ejercicio 1 usando scikit-learn. El código debe incluir:

\begin{enumerate}
    \item[a.] Preprocesamiento de datos categóricos
    \item[b.] División en entrenamiento y prueba
    \item[c.] Entrenamiento del modelo
    \item[d.] Evaluación con métricas apropiadas
    \item[e.] Matriz de confusión
\end{enumerate}

\begin{lstlisting}[language=Python]
# Codigo base para completar
import pandas as pd
from sklearn.naive_bayes import GaussianNB
from sklearn.preprocessing import LabelEncoder
from sklearn.model_selection import train_test_split
from sklearn.metrics import classification_report, confusion_matrix

# Completar la implementacion...
\end{lstlisting}

\subsection{Ejercicio 5 - Problema Médico}

Un médico quiere usar Naive Bayes para diagnosticar una enfermedad basándose en síntomas. Los datos históricos muestran:

\begin{itemize}
    \item P(Enfermedad) = 0.01 (1\% de la población tiene la enfermedad)
    \item P(Síntoma|Enfermedad) = 0.9 (90\% de los enfermos presenta el síntoma)
    \item P(Síntoma|Sano) = 0.05 (5\% de los sanos presenta el síntoma)
\end{itemize}

\begin{enumerate}
    \item[a.] Si un paciente presenta el síntoma, ¿cuál es la probabilidad de que tenga la enfermedad?
    \item[b.] ¿Por qué este resultado puede ser contraintuitivo?
    \item[c.] ¿Qué implicaciones tiene esto para el uso de Naive Bayes en medicina?
\end{enumerate}

\subsection{Ejercicio 6 - Clasificación de Texto Avanzada}

Para un problema de clasificación de texto, tienes un corpus con las siguientes características:

\begin{itemize}
    \item 10,000 documentos de entrenamiento
    \item 50,000 palabras únicas en el vocabulario
    \item 3 clases diferentes
    \item Algunos documentos muy cortos (2-3 palabras) y otros muy largos (1000+ palabras)
\end{itemize}

\begin{enumerate}
    \item[a.] ¿Qué variante de Naive Bayes usarías y por qué?
    \item[b.] ¿Cómo manejarías el problema de palabras que aparecen en test pero no en entrenamiento?
    \item[c.] ¿Aplicarías algún preprocesamiento adicional? ¿Cuál?
    \item[d.] ¿Cómo evaluarías el modelo considerando que las clases podrían estar desbalanceadas?
\end{enumerate}

\newpage

\section{Soluciones}

\subsection{Soluciones - Verdadero o Falso}

\begin{enumerate}
    \item \textbf{Falso.} Naive Bayes no requiere escalado de datos porque trabaja con probabilidades, no con distancias.
    \item \textbf{Verdadero.} Esta es la suposición fundamental del algoritmo.
    \item \textbf{Verdadero.} Es ideal para datos de conteo de palabras.
    \item \textbf{Verdadero.} Asume distribución gaussiana para características continuas.
    \item \textbf{Verdadero.} Evita el problema de probabilidades cero en datos de entrenamiento.
    \item \textbf{Falso.} El rendimiento depende del problema y los datos específicos.
    \item \textbf{Falso.} Aunque rara vez se cumple perfectamente, el algoritmo puede funcionar bien incluso con dependencias moderadas.
    \item \textbf{Verdadero.} Especialmente útil para características de presencia/ausencia.
    \item \textbf{Falso.} Es un algoritmo generativo, no discriminativo.
    \item \textbf{Falso.} Principalmente se usa para clasificación.
\end{enumerate}

\subsection{Soluciones - Opción Múltiple}

\begin{enumerate}
    \item \textbf{b)} Porque asume independencia entre características, lo cual rara vez es cierto en la realidad.
    \item \textbf{c)} Velocidad de entrenamiento y predicción - es muy eficiente computacionalmente.
    \item \textbf{b)} Multinomial - apropiado para conteos de palabras en documentos.
    \item \textbf{d)} Cuando hay probabilidades cero que pueden causar problemas en el cálculo.
    \item \textbf{a)} Puede reducir el rendimiento, pero el algoritmo sigue siendo robusto.
\end{enumerate}

\subsection{Soluciones - Ejercicios 1}

\begin{itemize}
    \item Contar observaciones por clase para calcular probabilidades a priori
    \item Para cada característica, calcular P(característica|clase) dividiendo conteos
    \item Para predecir: aplicar teorema de Bayes multiplicando probabilidades
    \item Accuracy final: 2 de 3 predicciones correctas
\end{itemize}

\subsection{Soluciones - Ejercicio 2}

\begin{itemize}
    \item P(Spam) = 4/6, P(No Spam) = 2/6
    \item Calcular probabilidades condicionales para cada palabra
    \item Cuidado con probabilidades cero en No Spam
    \item Con Laplace smoothing: agregar 1 al numerador y k al denominador
    \item El smoothing evita multiplicar por cero
\end{itemize}

\subsection{Soluciones - Ejercicio 3}

Bag of Words convierte texto en matriz numérica dispersa. Multinomial Naive Bayes es ideal para conteos de palabras en clasificación de documentos.

\subsection{Soluciones - Ejercicio 4}

\begin{itemize}
    \item Usar LabelEncoder para variables categóricas
    \item Aplicar train\_test\_split para división de datos
    \item Entrenar con CategoricalNB o MultinomialNB según el tipo de datos
    \item Evaluar con accuracy, precision, recall y matriz de confusión
\end{itemize}

\subsection{Soluciones - Ejercicio 5}

Aplicar teorema de Bayes: P(Enfermedad|Síntoma) $\approx$ 0.154 (15.4\%). La baja prevalencia (1\%) hace que sea menos probable tener la enfermedad incluso con síntoma positivo.

\subsection{Soluciones - Ejercicio 6}

Usar Multinomial Naive Bayes para texto. Aplicar Laplace smoothing para palabras nuevas. Preprocesar con tokenización, remoción de stopwords y stemming. Evaluar con F1-score macro para clases desbalanceadas.

\end{document}
